\subsection{Converting between Cartesian and polar descriptions}

The right triangle formed by the ray $OP$ gives
\[
\boxed{x=r\cos\theta,\qquad y=r\sin\theta}.
\]
Conversely,
\[
\boxed{r=\sqrt{x^2+y^2}}.
\]
The angle is determined by the direction of $(x,y)$; in computation this is the
role of the two-argument function $\operatorname{atan2}(y,x)$.

\begin{center}
\begin{tikzpicture}[scale=1.0,>=stealth]
    \draw[axis,->] (-0.8,0)--(6.0,0) node[right] {$x$};
    \draw[axis,->] (0,-0.8)--(0,4.8) node[above] {$y$};
    \coordinate (P) at (4.2,3.0);
    \draw[blue!70!black,line width=1.2pt,->] (0,0)--(P)
        node[midway,above left] {$r$};
    \draw[helper,dashed] (P)--(4.2,0);
    \draw[helper,dashed] (P)--(0,3.0);
    \draw (1.0,0) arc[start angle=0,end angle=35.54,radius=1.0];
    \node at (1.25,0.45) {$\theta$};
    \node at (2.1,-0.35) {$r\cos\theta$};
    \node[rotate=90] at (-0.4,1.5) {$r\sin\theta$};
    \fill[geometricSubsetColor] (P) circle (2.2pt) node[above right] {$P$};
\end{tikzpicture}
\end{center}

\begin{examplebox}
For $(r,\theta)=(4,\pi/3)$,
\[
x=4\cos\frac\pi3=2,
\qquad
y=4\sin\frac\pi3=2\sqrt3.
\]
Thus the same point has Cartesian coordinates $(2,2\sqrt3)$.
\end{examplebox}

\begin{corebox}[Coordinates are descriptions, not objects]
Changing from $(x,y)$ to $(r,\theta)$ does not move the point. It changes which
questions are easy to answer. Cartesian coordinates expose horizontal and
vertical displacement; polar coordinates expose direction and distance from a
chosen point.
\end{corebox}
